% หน้าอนุมัติและช่องลงนาม / Approval page and signature fields
\begingroup
\setlength{\parindent}{0pt}

% ตารางล่องหนทำให้เส้น ชื่อ และตำแหน่งตรงกัน / Invisible grid aligns signature lines, names, and roles
\newcommand{\signatureblock}[2]{%
  \noindent\begin{tabularx}{\textwidth}{@{}X>{\centering\arraybackslash}p{7cm}@{\hspace{0.5em}}p{3.3cm}@{}}
    & \hrulefill & #1 \\
    & (#2) & \\
  \end{tabularx}%
}

\begin{tabular}{@{}p{3.5cm}@{ : }p{10.1cm}@{}}
หัวข้องานวิจัย & \thesistitleth\\
               & \thesistitleen\\
ผู้วิจัย        & \authorname\\
รหัสประจำตัว   & \studentid\\
\ifsecondauthor
ผู้วิจัย        & \secondauthorname\\
รหัสประจำตัว   & \secondstudentid\\
\fi
สาขาวิชา       & วิทยาการคอมพิวเตอร์และปัญญาประดิษฐ์\\
อาจารย์ที่ปรึกษา & \advisor\\
\ifcoadvisor
อาจารย์ที่ปรึกษาร่วม & \coadvisor\\
\fi
\end{tabular}

\vspace{16pt}
\hrule
\vspace{16pt}

\setlength{\parindent}{1.27cm}
\degree\space \faculty\space \university\space
อนุมัติให้นับงานวิจัยนี้เป็นส่วนหนึ่งของการศึกษาตามหลักสูตรปริญญาบัณฑิต

\vspace{48pt}
\signatureblock{ประธานหลักสูตร}{\programchair}

\vspace{16pt}
\noindent คณะกรรมการสอบงานวิจัย

\vspace{32pt}
\signatureblock{ประธานกรรมการ}{\examchair}

\vspace{32pt}
\signatureblock{กรรมการ}{\committee}

\vspace{32pt}
\signatureblock{อาจารย์ที่ปรึกษา}{\advisor}

\ifcoadvisor
\vspace{32pt}
\signatureblock{อาจารย์ที่ปรึกษาร่วม}{\coadvisor}
\fi

\endgroup
